% Lösungen Zone (zusammenbau v0.1)
\einheitenkopf[][Kennst du schon]{Das kennst du schon}

Zuordnung: 1 → die besonderen Lagen in Einheit 1 und die Sachbilder in Einheit 4 · 2 → Richtungsvektoren in allen Einheiten, der Abstandspunkt in Einheit 2 · 3 → die senkrecht schneidende Gerade in Einheit 1 · 4 → der Parameter der Punktprobe in Einheit 2 · 5 → die Parallelitäts- und Identitätsfragen in Einheit 3 · 6 → die ebene Straßengerade in Einheit 4 und die Grundvorstellung des Fahrplans · 7 → Neigung und Realgrößen in Einheit 4 · 8 → der Parameter der Punktprobe in Einheit 2 · 9 → der Parameter der Punktprobe in Einheit 2

\erg{1}{a) in der $x_1x_3$-Ebene, weil die zweite Koordinate null ist \quad b) $(4 | -3 | -2{,}5)$}
\erg{2}{a) $\overrightarrow{AB} = (3 | 2 | 0)$ \quad b) $|\vec u| = 5$, also $(0 | 0{,}6 | -0{,}8)$}
\erg{3}{a) $2 + 4 + 6 = 12$ \quad b) $2 + 2a - 3 = 0$, also $a = 0{,}5$}
\erg{4}{a) $t = 3$ \quad b) $r = -1{,}5$}
\erg{5}{a) Ja, $\vec b = 3 \cdot \vec a$. \quad b) Ja: $(-2 | 6 | -8) = -4 \cdot (0{,}5 | -1{,}5 | 2)$.}
\erg{6}{a) $m = 3$, $n = -2$ \quad b) $n = 3$, also $y = -0{,}5x + 3$}
\erg{7}{a) $\frac{1}{20} = 5\,\%$ \quad b) $1\,250$ m $= 1{,}25$ km}
\erg{8}{a) Die dritte Gleichung wurde nicht geprüft: $-1 + 2 \cdot 3 = 5 \ne 4$ – es gibt keine solche Zahl.}
\erg{9}{a) Nein: aus der ersten folgt $s = 3$, die zweite stimmt, die dritte ergibt $7 \ne 8$.}
